% 只展示已有实现：冲突诊断不等同于已完成样本消解或降权。
\begingroup
\definecolor{qflowblue}{HTML}{315A79}
\definecolor{qflowlight}{HTML}{EDF3F7}
\definecolor{qflowgray}{HTML}{536572}
\begin{tikzpicture}[
  font=\song\fontsize{10}{13}\selectfont,
  qbox/.style={draw=qflowblue!55, line width=0.55pt, rounded corners=2pt,
    fill=qflowlight, align=center, inner sep=7pt, minimum height=1.15cm},
  qwide/.style={qbox, text width=6.0cm},
  qhalf/.style={qbox, text width=2.55cm, minimum height=1.7cm},
  qlink/.style={-{Stealth[length=2mm]}, draw=qflowblue, line width=0.7pt},
  qhead/.style={text=qflowblue, font=\song\bfseries\fontsize{11}{14}\selectfont}
]
\node[qhead] at (-3.9,0.9) {质量评价流程};
\node[qhead] at (3.9,0.9) {配比建模流程};
\node[qwide, fill=qflowblue, text=white] (quality) at (-3.9,0)
  {A1--A3 质量信号\\22 个指标与文本领域信息};
\node[qwide, fill=qflowblue, text=white] (mixture) at (3.9,0)
  {A4--A5 训练配方与验证损失\\17 个训练领域、13 个验证领域};

\node[qwide] (prepare) at (-3.9,-1.95)
  {缺失值处理与列表标量化\\统一指标方向并进行归一化};
\node[qwide] (match) at (3.9,-1.95)
  {按 index 匹配配方与损失\\构造配比矩阵 $\bm P$ 与损失矩阵 $\bm L$};

\node[qhalf] (score) at (-5.62,-4.15)
  {熵权综合评价\\样本、语料和领域\\质量评分};
\node[qhalf] (conflict) at (-2.18,-4.15)
  {指标冲突诊断\\成对标准化分数差\\Kendall 一致性};
\node[qwide] (ridge) at (3.9,-4.15)
  {岭正则线性混合回归\\估计回归系数并预测验证损失};

\node[qwide, fill=white] (mapping) at (-3.9,-6.55)
  {A16 领域映射\\构造混合质量指数 $\bar Q(\bm p)$};
\node[qwide] (validation) at (3.9,-6.55)
  {模型检验与跨尺度分析\\1M 检验、60M/1B 相关性\\10B/70B 外推表};
\node[qwide, fill=white] (recommend) at (3.9,-8.75)
  {有界配比推荐\\固定未测量域、设置配比上限};

\draw[qlink] (quality) -- (prepare);
\draw[qflowblue, line width=0.7pt] (prepare.south) -- (-3.9,-2.95);
\draw[qlink] (-3.9,-2.95) -| (score.north);
\draw[qlink] (-3.9,-2.95) -| (conflict.north);
\draw[qlink] (score.south) -- (-5.62,-5.7) -| (mapping.north);

\draw[qlink] (mixture) -- (match);
\draw[qlink] (match) -- (ridge);
\draw[qlink] (ridge) -- (validation);
\draw[qlink] (validation) -- (recommend);
\draw[qlink, dashed] (mapping.east) -- (0,-6.55) -- (0,-4.15) -- (ridge.west);
\node[fill=white, text=qflowgray, align=center, inner sep=2pt,
  font=\song\fontsize{8.5}{11}\selectfont] at (0,-5.35) {质量项\\增量检验};

\node[align=left, text=qflowgray, text width=6.0cm,
  font=\song\fontsize{9}{12}\selectfont] at (-3.9,-8.6)
  {实线表示主要处理流程；虚线表示混合质量指数的追加检验。};
\end{tikzpicture}
\endgroup
